%%%%%%%%%%%%%%%%%%%%%%
%%%%%%%%%%%%%%%%%%%%%%

\begin{figure*}
\begin{centering}
	\begin{subfigure}[b]{0.45\textwidth}
                \centering              
                \includegraphics[width=\textwidth]{Figures/oc/new_results/human_mouse_mfo_2}
                \caption{Human/Mouse}                  
                \label{fig:ocHMM2}
        \end{subfigure}%
	\quad
	\begin{subfigure}[b]{0.45\textwidth}
                \centering        
                \includegraphics[width=\textwidth]{Figures/oc/new_results/human_rat_mfo_2}
                \caption{Human/Rat}                        
                \label{fig:ocHRM2}
        \end{subfigure}%

	\begin{subfigure}[b]{0.45\textwidth}
                \centering       
                \includegraphics[width=\textwidth]{Figures/oc/new_results/mouse_rat_mfo_2}
                \caption{Mouse/Rat}                         
                \label{fig:ocMRM2}
        \end{subfigure}%
	\quad
	\begin{subfigure}[b]{0.45\textwidth}
                \centering     
                \includegraphics[width=\textwidth]{Figures/oc/new_results/worm_fly_mfo_2}
                \caption{Worm/Fly}                           
                \label{fig:ocWFM2}
        \end{subfigure}%

\end{centering}

\caption[Comparison of orthologs and paralogs in four pairs of organisms for Molecular Function annotations]{Comparison of orthologs (red) and paralog (blue) functional similarity across a range of global sequence identity values in four pairs of organisms when considering Molecular Function annotations. Standard error bars are shown. }
\label{fig:oc1_mfo}
\end{figure*}

%%%%%%%%%%%%%%%%%%%%%%
%%%%%%%%%%%%%%%%%%%%%%